\documentclass[10pt,a4paper]{amsart}
\usepackage{amsmath,amssymb,amsthm,graphicx}
\usepackage{verbatim}
\newenvironment{reference}{\comment}{\endcomment}
%\newenvironment{reference}{}{}
\newenvironment{slogan}{\comment}{\endcomment}
\newenvironment{history}{\comment}{\endcomment}

% For commutative diagrams we use Xy-pic
\usepackage[all]{xy}

% We use 2cell for 2-commutative diagrams.
\xyoption{2cell}
% Theorem environments.
%
\theoremstyle{plain}
\newtheorem{theorem}[subsection]{Theorem}
\newtheorem{proposition}[subsection]{Proposition}
\newtheorem{lemma}[subsection]{Lemma}

\theoremstyle{definition}
\newtheorem{definition}[subsection]{Definition}
\newtheorem{example}[subsection]{Example}
\newtheorem{exercise}[subsection]{Exercise}
\newtheorem{situation}[subsection]{Situation}

\theoremstyle{remark}
\newtheorem{remark}[subsection]{Remark}
\newtheorem{remarks}[subsection]{Remarks}

\numberwithin{equation}{subsection}

% Macros
%
\def\lim{\mathop{\mathrm{lim}}\nolimits}
\def\colim{\mathop{\mathrm{colim}}\nolimits}
\def\Spec{\mathop{\mathrm{Spec}}}
\def\Hom{\mathop{\mathrm{Hom}}\nolimits}
\def\Ext{\mathop{\mathrm{Ext}}\nolimits}
\def\SheafHom{\mathop{\mathcal{H}\!\mathit{om}}\nolimits}
\def\SheafExt{\mathop{\mathcal{E}\!\mathit{xt}}\nolimits}
\def\Sch{\mathit{Sch}}
\def\Mor{\mathop{\mathrm{Mor}}\nolimits}
\def\Ob{\mathop{\mathrm{Ob}}\nolimits}
\def\Sh{\mathop{\mathit{Sh}}\nolimits}
\def\NL{\mathop{N\!L}\nolimits}
\def\CH{\mathop{\mathrm{CH}}\nolimits}
\def\proetale{{pro\text{-}\acute{e}tale}}
\def\etale{{\acute{e}tale}}
\def\QCoh{\mathit{QCoh}}
\def\Ker{\mathop{\mathrm{Ker}}}
\def\Im{\mathop{\mathrm{Im}}}
\def\Coker{\mathop{\mathrm{Coker}}}
\def\Coim{\mathop{\mathrm{Coim}}}

% Boxtimes
%
\DeclareMathSymbol{\boxtimes}{\mathbin}{AMSa}{"02}

%
% Macros for moduli stacks/spaces
%
\def\QCohstack{\mathcal{QC}\!\mathit{oh}}
\def\Cohstack{\mathcal{C}\!\mathit{oh}}
\def\Spacesstack{\mathcal{S}\!\mathit{paces}}
\def\Quotfunctor{\mathrm{Quot}}
\def\Hilbfunctor{\mathrm{Hilb}}
\def\Curvesstack{\mathcal{C}\!\mathit{urves}}
\def\Polarizedstack{\mathcal{P}\!\mathit{olarized}}
\def\Complexesstack{\mathcal{C}\!\mathit{omplexes}}
% \Pic is the operator that assigns to X its picard group, usage \Pic(X)
% \Picardstack_{X/B} denotes the Picard stack of X over B
% \Picardfunctor_{X/B} denotes the Picard functor of X over B
\def\Pic{\mathop{\mathrm{Pic}}\nolimits}
\def\Picardstack{\mathcal{P}\!\mathit{ic}}
\def\Picardfunctor{\mathrm{Pic}}
\def\Deformationcategory{\mathcal{D}\!\mathit{ef}}
\usepackage{cleveref}
\crefname{originalreference}{original reference}{original references}
\Crefname{originalreference}{Original reference}{Original references}
\makeatletter
\newlabel{remark-compare-S1-S2-schlessinger}{{06HY}{source}{}{sourcepointer-1}{}}
\newlabel{remark-compare-S1-S2-schlessinger@cref}{{[originalreference][0][]06HY}{[1][0][]source}{}{}{}}
\makeatother
\title{Condition (S2) with the trivial extension replaced by an arbitrary finite dimensional coefficient space}
\author{The Stacks Project Authors}
\date{stacks-project:formal-defos.tex:local-snapshot-20261012; GNU FDL 1.2 or later}
\begin{document}
\maketitle
\noindent\textit{Editorial note.} All delivered mathematical text is the author's own, copied from the pinned Stacks Project chapter file formal-defos.tex; no statement or proof step was rewritten, reordered or removed.\par\medskip
\noindent\textit{Editorial note.} The counted proof is the author's complete argument: the case of the associated category in sets through the Schlessinger description, the induction on the dimension with the splitting V = V' direct sum k epsilon, and the separate treatments of the existence and injectivity halves using the two parts of (S2).\par\medskip
\noindent\textit{Editorial note.} The label lemma-S2-extensions names the counted statement, and the proof environment immediately following it in the source is delivered verbatim as its proof.\par\medskip
\noindent\textit{Editorial note.} The counted statement is carried by the author's own lemma environment in the chapter source.\par\medskip
\noindent\textit{Editorial note.} Pre-open gate: python operations/agent-authoring/tools/precard\_gate.py --source formal-defos.tex --label lemma-S2-extensions --start 2871 --end 3092 --local-id 28030 returned RESULT=OK (L28030 free, occupancy FREE, public interval clear). Candidate taken from sources/stacks-formal-defos/candidates-halfA.json.\par\medskip
\noindent\textit{Editorial note.} The statement and proof are copied from the pinned chapter file; the cited inputs are kept as external pointers with their canonical Stacks tags.\par\medskip
\noindent\textit{Allowed standard prerequisites (editorial).}\begin{itemize}
\item Description of (S2) for a functor in terms of fibre products over k. Same chapter, the remark comparing the categorical formulation of (S1) and (S2) with the Schlessinger formulation for functors. Supplies the case of the associated category in sets, where k[V] is a product of copies of k[epsilon].
\item Split exactness of the coefficient algebra and the identification of the fibre products. Standard algebra identifications k[V] = k[V'] x\_k k[epsilon] for a decomposition V = V' direct sum k epsilon, and the compatibility of the ring with fibre products over k. Supplies the identity (A x\_k k[V']) x\_k k[epsilon] = A x\_k k[V] used at the key step.
\item The two halves of condition (S2). Same chapter, the statement of (S2) as the existence of objects over A x\_k k[epsilon] plus the injectivity of the comparison map. Supplies both the existence step and the injectivity step applied along the induction.
\end{itemize}\medskip
\ifdefined\hypertarget\hypertarget{sourcepointer-1}{}\fi
\noindent\textit{Source reference.} Original reference 06HY: sources/stacks-formal-defos/pinned-original-source.tex, source lines 2850--2850. Prerequisite pointer invoked at proof line(s) 2883; it names a standard fact of the same chapter that the counted argument uses, and every step of the counted proof built on it is retained verbatim in the delivered body.\par\medskip
\begin{lemma}
\label{lemma-S2-extensions}
Let $\mathcal{F}$ be a category cofibred in groupoids over
$\mathcal{C}_\Lambda$. If $\mathcal{F}$ satisfies (S2), then the
condition of (S2) also holds when $k[\epsilon]$ is replaced by $k[V]$
for any finite dimensional $k$-vector space $V$.
\end{lemma}

\begin{proof}
In the case that $\mathcal{F}$ is cofibred in sets, i.e., corresponds
to a functor $F : \mathcal{C}_\Lambda \to \textit{Sets}$ this follows
from the description of (S2) for $F$ in
Remark \ref{remark-compare-S1-S2-schlessinger}
and the fact that
$k[V] \cong k[\epsilon] \times_k \ldots \times_k k[\epsilon]$
with $\dim_k V$ factors. The case of functors is what we will use in
the rest of this chapter.

\medskip\noindent
We prove the general case by induction on $\dim(V)$. If $\dim(V) = 1$, then
$k[V] \cong k[\epsilon]$ and the result holds by assumption.
If $\dim(V) > 1$ we write $V = V' \oplus k\epsilon$. Pick a diagram
$$
\vcenter{
\xymatrix{
& x_V \ar[d] \\
x \ar[r] & x_0
}
}
\quad\text{lying over}\quad
\vcenter{
\xymatrix{
& k[V] \ar[d] \\
A \ar[r] & k
}
}
$$
Choose a morphism $x_V \to x_{V'}$ lying over $k[V] \to k[V']$
and a morphism $x_V \to x_\epsilon$ lying over $k[V] \to k[\epsilon]$.
Note that the morphism $x_V \to x_0$ factors as
$x_V \to x_{V'} \to x_0$ and as $x_V \to x_\epsilon \to x_0$.
By induction hypothesis we can find a diagram
$$
\vcenter{
\xymatrix{
y' \ar[d] \ar[r] & x_{V'} \ar[d] \\
x \ar[r] & x_0
}
}
\quad\text{lying over}\quad
\vcenter{
\xymatrix{
A \times_k k[V'] \ar[d] \ar[r] & k[V'] \ar[d] \\
A \ar[r] & k
}
}
$$
This gives us a commutative diagram
$$
\vcenter{
\xymatrix{
& x_\epsilon \ar[d] \\
y' \ar[r] & x_0
}
}
\quad\text{lying over}\quad
\vcenter{
\xymatrix{
& k[\epsilon] \ar[d] \\
A \times_k k[V'] \ar[r] & k
}
}
$$
Hence by (S2) we get a commutative diagram
$$
\vcenter{
\xymatrix{
y \ar[d] \ar[r] & x_\epsilon \ar[d] \\
y' \ar[r] & x_0
}
}
\quad\text{lying over}\quad
\vcenter{
\xymatrix{
(A \times_k k[V']) \times_k k[\epsilon] \ar[d] \ar[r] & k[\epsilon] \ar[d] \\
A \times_k k[V'] \ar[r] & k
}
}
$$
Note that
$(A \times_k k[V']) \times_k k[\epsilon] = A \times_k k[V' \oplus k\epsilon]
= A \times_k k[V]$. We claim that $y$ fits into the correct commutative
diagram. To see this we let $y \to y_V$ be a morphism lying over
$A \times_k k[V] \to k[V]$. We can factor the morphisms
$y \to y' \to x_{V'}$ and $y \to x_\epsilon$ through the morphism
$y \to y_V$ (by the axioms of categories cofibred in groupoids).
Hence we see that both $y_V$ and $x_V$ fit into commutative diagrams
$$
\vcenter{
\xymatrix{
y_V \ar[r] \ar[d] & x_\epsilon \ar[d] \\
x_{V'} \ar[r]          & x_0
}
}
\quad\text{and}\quad
\vcenter{
\xymatrix{
x_V \ar[r] \ar[d] & x_\epsilon \ar[d] \\
x_{V'} \ar[r]                 & x_0
}
}
$$
and hence by the second part of (S2) there exists an isomorphism
$y_V \to x_V$ compatible with $y_V \to x_{V'}$ and $x_V \to x_{V'}$
and in particular compatible with the maps to $x_0$.
The composition $y \to y_V \to x_V$ then fits into the required commutative
diagram
$$
\vcenter{
\xymatrix{
y \ar[r] \ar[d] & x_V \ar[d] \\
x \ar[r] & x_0
}
}
\quad\text{lying over}\quad
\vcenter{
\xymatrix{
A \times_k k[V] \ar[d] \ar[r] & k[V] \ar[d] \\
A \ar[r] & k
}
}
$$
In this way we see that the first part of $(S2)$ holds with $k[\epsilon]$
replaced by $k[V]$.

\medskip\noindent
To prove the second part suppose given two commutative
diagrams
$$
\vcenter{
\xymatrix{
y \ar[r] \ar[d] & x_V \ar[d] \\
x \ar[r] & x_0
}
}
\quad\text{and}\quad
\vcenter{
\xymatrix{
y' \ar[r] \ar[d] & x_V \ar[d] \\
x \ar[r] & x_0
}
}
\quad\text{lying over}\quad
\vcenter{
\xymatrix{
A \times_k k[V] \ar[d] \ar[r] & k[V] \ar[d] \\
A \ar[r] & k
}
}
$$
We will use the morphisms $x_V \to x_{V'} \to x_0$ and
$x_V \to x_\epsilon \to x_0$ introduced in the first paragraph of the proof.
Choose morphisms $y \to y_{V'}$ and $y' \to y'_{V'}$
lying over $A \times_k k[V] \to A \times_k k[V']$. The axioms of a
cofibred category imply we can find commutative diagrams
$$
\vcenter{
\xymatrix{
y_{V'} \ar[r] \ar[d] & x_{V'} \ar[d] \\
x \ar[r] & x_0
}
}
\quad\text{and}\quad
\vcenter{
\xymatrix{
y'_{V'} \ar[r] \ar[d] & x_{V'} \ar[d] \\
x \ar[r] & x_0
}
}
\quad\text{lying over}\quad
\vcenter{
\xymatrix{
A \times_k k[V'] \ar[d] \ar[r] & k[V'] \ar[d] \\
A \ar[r] & k
}
}
$$
By induction hypothesis we obtain an isomorphism
$b : y_{V'} \to y'_{V'}$
compatible with the morphisms $y_{V'} \to x$ and $y'_{V'} \to x$,
in particular compatible with the morphisms to $x_0$.
Then we have commutative diagrams
$$
\vcenter{
\xymatrix{
y \ar[r] \ar[d] & x_\epsilon \ar[d] \\
y'_{V'} \ar[r] & x_0
}
}
\quad\text{and}\quad
\vcenter{
\xymatrix{
y' \ar[r] \ar[d] & x_\epsilon \ar[d] \\
y'_{V'} \ar[r] & x_0
}
}
\quad\text{lying over}\quad
\vcenter{
\xymatrix{
A \times_k k[\epsilon] \ar[d] \ar[r] & k[\epsilon] \ar[d] \\
A \ar[r] & k
}
}
$$
where the morphism $y \to y'_{V'}$ is the composition
$y \to y_{V'} \xrightarrow{b} y'_{V'}$ and where the morphisms
$y \to x_\epsilon$ and $y' \to x_\epsilon$ are the compositions of
the maps $y \to x_V$ and $y' \to x_V$ with the morphism $x_V \to x_\epsilon$.
Then the second part of (S2) guarantees the existence of an isomorphism
$y \to y'$ compatible with the maps to $y'_{V'}$, in particular compatible
with the maps to $x$ (because $b$ was compatible with the maps to $x$).
\end{proof}

\end{document}
